% =====================================================================
% P5 merge snippet — M02: evidence protocol (six rules + its own limits)
% Source: manuscript_v2, Sec. 2.7 (sec:evidence_protocol), written 2026-11-01.
% Adapted: \cref/\eqref, this paper's appendix labels, \citep.
% NEW LABELS: sec:evidence_protocol
% Suggested placement: a new subsection at the end of \cref{sec:intro_ref} (or after the
% reporting standards in \cref{sec:reporting}), because it tells the reader how to discount
% every statement in between.
% =====================================================================

\subsection{Evidence protocol and how to discount this review}\label{sec:evidence_protocol}

A review is only as useful as the reader's ability to tell what it has checked from what it has merely
repeated. The protocol applied throughout is stated here so that any statement can be discounted by the
rule it fails.

\begin{enumerate}
\item \textbf{Derivation and attribution are separated.} A statement derived here is presented as a
derivation and carries an equation identifier; a statement taken from a source is presented as
attribution even where the source's formulation is reproduced exactly. The two are never merged into one
unmarked sentence. The register in \cref{sec:eqregister} records which is which.
\item \textbf{Every registered equation carries a status} from a fixed vocabulary, and a status records
the check performed, not a rank of importance (\cref{tab:eqregister}).
\item \textbf{A secondary source is never promoted to a primary one.} Where the primary text was not
available the text says so and the register records the weaker status. The absence of the primary text
is not evidence for the claim.
\item \textbf{Bibliographic identifiers are verified before they are printed.} A DOI printed in the
reference list has been checked against a bibliographic registry record for the same work --- author,
title, journal, volume, pages --- rather than inferred from a citation in another paper. Works that
predate DOI assignment say so in the entry instead of carrying a plausible-looking identifier.
\item \textbf{Numbers that can be recomputed are given with what is needed to recompute them.} A benchmark
value is presented with its specification, reference solution and tolerance
(\cref{app:registry}); a value quoted from the literature is presented as quoted, with its conditions,
and not as an independent confirmation.
\item \textbf{Errors are recorded without personal characterisation.} Where a material error is
conclusively established it is recorded plainly; where the evidence establishes a difference but not an
error, it is recorded as a difference. No allegation is made against a publication, and no model-specific
discrepancy is presented as a mathematical contradiction without demonstration.
\end{enumerate}

The protocol has limits and they are part of it. Not every cited work was read page by page; coverage of
the literature is deliberate rather than exhaustive; and the register documents checks performed for this
version, not checks that exist somewhere in principle. Those limits are why the register reports coverage
instead of asserting completeness.
